% Options for packages loaded elsewhere
\PassOptionsToPackage{unicode}{hyperref}
\PassOptionsToPackage{hyphens}{url}
\documentclass[
]{article}
\usepackage{xcolor}
\usepackage[margin=1in]{geometry}
\usepackage{amsmath,amssymb}
\setcounter{secnumdepth}{-\maxdimen} % remove section numbering
\usepackage{iftex}
\ifPDFTeX
  \usepackage[T1]{fontenc}
  \usepackage[utf8]{inputenc}
  \usepackage{textcomp} % provide euro and other symbols
\else % if luatex or xetex
  \usepackage{unicode-math} % this also loads fontspec
  \defaultfontfeatures{Scale=MatchLowercase}
  \defaultfontfeatures[\rmfamily]{Ligatures=TeX,Scale=1}
\fi
\usepackage{lmodern}
\ifPDFTeX\else
  % xetex/luatex font selection
\fi
% Use upquote if available, for straight quotes in verbatim environments
\IfFileExists{upquote.sty}{\usepackage{upquote}}{}
\IfFileExists{microtype.sty}{% use microtype if available
  \usepackage[]{microtype}
  \UseMicrotypeSet[protrusion]{basicmath} % disable protrusion for tt fonts
}{}
\makeatletter
\@ifundefined{KOMAClassName}{% if non-KOMA class
  \IfFileExists{parskip.sty}{%
    \usepackage{parskip}
  }{% else
    \setlength{\parindent}{0pt}
    \setlength{\parskip}{6pt plus 2pt minus 1pt}}
}{% if KOMA class
  \KOMAoptions{parskip=half}}
\makeatother
\usepackage{color}
\usepackage{fancyvrb}
\newcommand{\VerbBar}{|}
\newcommand{\VERB}{\Verb[commandchars=\\\{\}]}
\DefineVerbatimEnvironment{Highlighting}{Verbatim}{commandchars=\\\{\}}
% Add ',fontsize=\small' for more characters per line
\usepackage{framed}
\definecolor{shadecolor}{RGB}{248,248,248}
\newenvironment{Shaded}{\begin{snugshade}}{\end{snugshade}}
\newcommand{\AlertTok}[1]{\textcolor[rgb]{0.94,0.16,0.16}{#1}}
\newcommand{\AnnotationTok}[1]{\textcolor[rgb]{0.56,0.35,0.01}{\textbf{\textit{#1}}}}
\newcommand{\AttributeTok}[1]{\textcolor[rgb]{0.13,0.29,0.53}{#1}}
\newcommand{\BaseNTok}[1]{\textcolor[rgb]{0.00,0.00,0.81}{#1}}
\newcommand{\BuiltInTok}[1]{#1}
\newcommand{\CharTok}[1]{\textcolor[rgb]{0.31,0.60,0.02}{#1}}
\newcommand{\CommentTok}[1]{\textcolor[rgb]{0.56,0.35,0.01}{\textit{#1}}}
\newcommand{\CommentVarTok}[1]{\textcolor[rgb]{0.56,0.35,0.01}{\textbf{\textit{#1}}}}
\newcommand{\ConstantTok}[1]{\textcolor[rgb]{0.56,0.35,0.01}{#1}}
\newcommand{\ControlFlowTok}[1]{\textcolor[rgb]{0.13,0.29,0.53}{\textbf{#1}}}
\newcommand{\DataTypeTok}[1]{\textcolor[rgb]{0.13,0.29,0.53}{#1}}
\newcommand{\DecValTok}[1]{\textcolor[rgb]{0.00,0.00,0.81}{#1}}
\newcommand{\DocumentationTok}[1]{\textcolor[rgb]{0.56,0.35,0.01}{\textbf{\textit{#1}}}}
\newcommand{\ErrorTok}[1]{\textcolor[rgb]{0.64,0.00,0.00}{\textbf{#1}}}
\newcommand{\ExtensionTok}[1]{#1}
\newcommand{\FloatTok}[1]{\textcolor[rgb]{0.00,0.00,0.81}{#1}}
\newcommand{\FunctionTok}[1]{\textcolor[rgb]{0.13,0.29,0.53}{\textbf{#1}}}
\newcommand{\ImportTok}[1]{#1}
\newcommand{\InformationTok}[1]{\textcolor[rgb]{0.56,0.35,0.01}{\textbf{\textit{#1}}}}
\newcommand{\KeywordTok}[1]{\textcolor[rgb]{0.13,0.29,0.53}{\textbf{#1}}}
\newcommand{\NormalTok}[1]{#1}
\newcommand{\OperatorTok}[1]{\textcolor[rgb]{0.81,0.36,0.00}{\textbf{#1}}}
\newcommand{\OtherTok}[1]{\textcolor[rgb]{0.56,0.35,0.01}{#1}}
\newcommand{\PreprocessorTok}[1]{\textcolor[rgb]{0.56,0.35,0.01}{\textit{#1}}}
\newcommand{\RegionMarkerTok}[1]{#1}
\newcommand{\SpecialCharTok}[1]{\textcolor[rgb]{0.81,0.36,0.00}{\textbf{#1}}}
\newcommand{\SpecialStringTok}[1]{\textcolor[rgb]{0.31,0.60,0.02}{#1}}
\newcommand{\StringTok}[1]{\textcolor[rgb]{0.31,0.60,0.02}{#1}}
\newcommand{\VariableTok}[1]{\textcolor[rgb]{0.00,0.00,0.00}{#1}}
\newcommand{\VerbatimStringTok}[1]{\textcolor[rgb]{0.31,0.60,0.02}{#1}}
\newcommand{\WarningTok}[1]{\textcolor[rgb]{0.56,0.35,0.01}{\textbf{\textit{#1}}}}
\usepackage{graphicx}
\makeatletter
\newsavebox\pandoc@box
\newcommand*\pandocbounded[1]{% scales image to fit in text height/width
  \sbox\pandoc@box{#1}%
  \Gscale@div\@tempa{\textheight}{\dimexpr\ht\pandoc@box+\dp\pandoc@box\relax}%
  \Gscale@div\@tempb{\linewidth}{\wd\pandoc@box}%
  \ifdim\@tempb\p@<\@tempa\p@\let\@tempa\@tempb\fi% select the smaller of both
  \ifdim\@tempa\p@<\p@\scalebox{\@tempa}{\usebox\pandoc@box}%
  \else\usebox{\pandoc@box}%
  \fi%
}
% Set default figure placement to htbp
\def\fps@figure{htbp}
\makeatother
\setlength{\emergencystretch}{3em} % prevent overfull lines
\providecommand{\tightlist}{%
  \setlength{\itemsep}{0pt}\setlength{\parskip}{0pt}}
\usepackage{bookmark}
\IfFileExists{xurl.sty}{\usepackage{xurl}}{} % add URL line breaks if available
\urlstyle{same}
\hypersetup{
  pdftitle={RNA-seq with DESeq2 / edgeR},
  hidelinks,
  pdfcreator={LaTeX via pandoc}}

\title{RNA-seq with DESeq2 / edgeR}
\author{}
\date{\vspace{-2.5em}}

\begin{document}
\maketitle

\textbf{Workflow labs} use the variant template: Goal → Approach →
Execution → Check → Report.

\subsection{Learning objectives}\label{learning-objectives}

\begin{itemize}
\tightlist
\item
  Describe the count-model assumptions behind DESeq2 and edgeR and
  explain the role of dispersion estimation.
\item
  Set up a design matrix for a two-group comparison and for a more
  complex experiment.
\item
  Recognise when a log-transformed \texttt{limma-voom} workflow is
  preferable.
\end{itemize}

\subsection{Prerequisites}\label{prerequisites}

GLMs; multiple testing.

\subsection{Background}\label{background}

Bulk RNA-sequencing yields integer counts per gene per sample. Raw
counts are heteroscedastic --- variance grows with the mean --- and the
distribution has a long tail. DESeq2 and edgeR model counts as
negative-binomial, with a gene-specific dispersion shrunk toward a trend
estimated across the dataset. This shrinkage is the statistical heart of
modern differential-expression analysis in low-replicate studies.

\texttt{limma-voom} takes a different route: model the log-counts with
precision weights that reflect the count-level mean--variance trend. The
output is a familiar linear-model workflow with all of its tools
(contrasts, interactions, F-tests). For large or complex designs,
\texttt{limma-voom} is often the most flexible.

DESeq2 and edgeR are Bioconductor packages, not CRAN. We show their call
pattern with \texttt{\#\textbar{}\ eval:\ false} and walk through a
conceptually parallel simulation with base R GLMs so the ideas run
end-to-end in any environment.

\subsection{Setup}\label{setup}

\begin{Shaded}
\begin{Highlighting}[]
\FunctionTok{library}\NormalTok{(tidyverse)}
\FunctionTok{library}\NormalTok{(MASS)}
\FunctionTok{set.seed}\NormalTok{(}\DecValTok{42}\NormalTok{)}
\FunctionTok{theme\_set}\NormalTok{(}\FunctionTok{theme\_minimal}\NormalTok{(}\AttributeTok{base\_size =} \DecValTok{12}\NormalTok{))}
\end{Highlighting}
\end{Shaded}

\subsection{1. Goal}\label{goal}

Simulate a small RNA-seq count matrix with two conditions and
demonstrate the DESeq2/edgeR pipeline pattern alongside a base-R
negative-binomial GLM that actually runs.

\subsection{2. Approach}\label{approach}

Simulate 500 genes × 10 samples (5 per group). Ten genes are truly
differentially expressed.

\begin{Shaded}
\begin{Highlighting}[]
\NormalTok{group }\OtherTok{\textless{}{-}} \FunctionTok{rep}\NormalTok{(}\FunctionTok{c}\NormalTok{(}\StringTok{"A"}\NormalTok{, }\StringTok{"B"}\NormalTok{), }\AttributeTok{each =} \DecValTok{5}\NormalTok{)}
\NormalTok{mu }\OtherTok{\textless{}{-}} \DecValTok{2} \SpecialCharTok{*} \FunctionTok{runif}\NormalTok{(n\_gene, }\DecValTok{1}\NormalTok{, }\DecValTok{10}\NormalTok{)   }\CommentTok{\# baseline mean per gene}
\NormalTok{lfc }\OtherTok{\textless{}{-}} \FunctionTok{c}\NormalTok{(}\FunctionTok{rep}\NormalTok{(}\FunctionTok{log}\NormalTok{(}\DecValTok{3}\NormalTok{), }\DecValTok{10}\NormalTok{), }\FunctionTok{rep}\NormalTok{(}\DecValTok{0}\NormalTok{, n\_gene }\SpecialCharTok{{-}} \DecValTok{10}\NormalTok{))}
\NormalTok{counts }\OtherTok{\textless{}{-}} \FunctionTok{sapply}\NormalTok{(}\FunctionTok{seq\_len}\NormalTok{(n\_samp), }\ControlFlowTok{function}\NormalTok{(j) \{}
\NormalTok{  rate }\OtherTok{\textless{}{-}}\NormalTok{ mu }\SpecialCharTok{*} \FunctionTok{exp}\NormalTok{(}\FunctionTok{ifelse}\NormalTok{(group[j] }\SpecialCharTok{==} \StringTok{"B"}\NormalTok{, lfc, }\DecValTok{0}\NormalTok{))}
  \FunctionTok{rnbinom}\NormalTok{(n\_gene, }\AttributeTok{mu =}\NormalTok{ rate, }\AttributeTok{size =} \DecValTok{5}\NormalTok{)}
\NormalTok{\})}
\FunctionTok{rownames}\NormalTok{(counts) }\OtherTok{\textless{}{-}} \FunctionTok{paste0}\NormalTok{(}\StringTok{"g"}\NormalTok{, }\FunctionTok{seq\_len}\NormalTok{(n\_gene))}
\FunctionTok{colnames}\NormalTok{(counts) }\OtherTok{\textless{}{-}} \FunctionTok{paste0}\NormalTok{(}\StringTok{"s"}\NormalTok{, }\FunctionTok{seq\_len}\NormalTok{(n\_samp))}

\FunctionTok{tibble}\NormalTok{(}\AttributeTok{mean =} \FunctionTok{rowMeans}\NormalTok{(counts), }\AttributeTok{var =} \FunctionTok{apply}\NormalTok{(counts, }\DecValTok{1}\NormalTok{, var)) }\SpecialCharTok{|\textgreater{}}
  \FunctionTok{ggplot}\NormalTok{(}\FunctionTok{aes}\NormalTok{(mean, var)) }\SpecialCharTok{+} \FunctionTok{geom\_point}\NormalTok{(}\AttributeTok{alpha =} \FloatTok{0.3}\NormalTok{) }\SpecialCharTok{+}
  \FunctionTok{scale\_x\_log10}\NormalTok{() }\SpecialCharTok{+} \FunctionTok{scale\_y\_log10}\NormalTok{() }\SpecialCharTok{+}
  \FunctionTok{geom\_abline}\NormalTok{(}\AttributeTok{slope =} \DecValTok{1}\NormalTok{, }\AttributeTok{intercept =} \DecValTok{0}\NormalTok{, }\AttributeTok{colour =} \StringTok{"grey50"}\NormalTok{) }\SpecialCharTok{+}
  \FunctionTok{labs}\NormalTok{(}\AttributeTok{title =} \StringTok{"mean–variance (log–log)"}\NormalTok{)}
\end{Highlighting}
\end{Shaded}

\subsection{3. Execution}\label{execution}

A DESeq2 call pattern (not run).

\begin{Shaded}
\begin{Highlighting}[]
\FunctionTok{library}\NormalTok{(DESeq2)}
\NormalTok{dds }\OtherTok{\textless{}{-}} \FunctionTok{DESeqDataSetFromMatrix}\NormalTok{(}
  \AttributeTok{countData =}\NormalTok{ counts,}
  \AttributeTok{colData   =} \FunctionTok{data.frame}\NormalTok{(}\AttributeTok{group =}\NormalTok{ group),}
  \AttributeTok{design    =} \SpecialCharTok{\textasciitilde{}}\NormalTok{ group}
\NormalTok{)}
\NormalTok{dds }\OtherTok{\textless{}{-}} \FunctionTok{DESeq}\NormalTok{(dds)}
\NormalTok{res }\OtherTok{\textless{}{-}} \FunctionTok{results}\NormalTok{(dds, }\AttributeTok{contrast =} \FunctionTok{c}\NormalTok{(}\StringTok{"group"}\NormalTok{, }\StringTok{"B"}\NormalTok{, }\StringTok{"A"}\NormalTok{))}
\FunctionTok{head}\NormalTok{(res[}\FunctionTok{order}\NormalTok{(res}\SpecialCharTok{$}\NormalTok{padj), ])}
\end{Highlighting}
\end{Shaded}

An edgeR call pattern (not run).

\begin{Shaded}
\begin{Highlighting}[]
\FunctionTok{library}\NormalTok{(edgeR)}
\NormalTok{dge }\OtherTok{\textless{}{-}} \FunctionTok{DGEList}\NormalTok{(}\AttributeTok{counts =}\NormalTok{ counts, }\AttributeTok{group =}\NormalTok{ group)}
\NormalTok{dge }\OtherTok{\textless{}{-}} \FunctionTok{calcNormFactors}\NormalTok{(dge)}
\NormalTok{design }\OtherTok{\textless{}{-}} \FunctionTok{model.matrix}\NormalTok{(}\SpecialCharTok{\textasciitilde{}}\NormalTok{ group)}
\NormalTok{dge }\OtherTok{\textless{}{-}} \FunctionTok{estimateDisp}\NormalTok{(dge, design)}
\NormalTok{fit }\OtherTok{\textless{}{-}} \FunctionTok{glmQLFit}\NormalTok{(dge, design)}
\NormalTok{qlf }\OtherTok{\textless{}{-}} \FunctionTok{glmQLFTest}\NormalTok{(fit, }\AttributeTok{coef =} \DecValTok{2}\NormalTok{)}
\FunctionTok{topTags}\NormalTok{(qlf)}
\end{Highlighting}
\end{Shaded}

A per-gene negative-binomial GLM that runs.

\begin{Shaded}
\begin{Highlighting}[]
\NormalTok{logfcs }\OtherTok{\textless{}{-}} \FunctionTok{numeric}\NormalTok{(n\_gene)}
\ControlFlowTok{for}\NormalTok{ (g }\ControlFlowTok{in} \FunctionTok{seq\_len}\NormalTok{(n\_gene)) \{}
\NormalTok{  y }\OtherTok{\textless{}{-}}\NormalTok{ counts[g, ]}
\NormalTok{  fit }\OtherTok{\textless{}{-}} \FunctionTok{try}\NormalTok{(}\FunctionTok{glm.nb}\NormalTok{(y }\SpecialCharTok{\textasciitilde{}}\NormalTok{ group), }\AttributeTok{silent =} \ConstantTok{TRUE}\NormalTok{)}
  \ControlFlowTok{if}\NormalTok{ (}\FunctionTok{inherits}\NormalTok{(fit, }\StringTok{"try{-}error"}\NormalTok{)) \{ pvals[g] }\OtherTok{\textless{}{-}} \ConstantTok{NA}\NormalTok{; }\ControlFlowTok{next}\NormalTok{ \}}
\NormalTok{  s }\OtherTok{\textless{}{-}} \FunctionTok{summary}\NormalTok{(fit)}
\NormalTok{  pvals[g]  }\OtherTok{\textless{}{-}}\NormalTok{ s}\SpecialCharTok{$}\NormalTok{coefficients[}\DecValTok{2}\NormalTok{, }\DecValTok{4}\NormalTok{]}
\NormalTok{  logfcs[g] }\OtherTok{\textless{}{-}}\NormalTok{ s}\SpecialCharTok{$}\NormalTok{coefficients[}\DecValTok{2}\NormalTok{, }\DecValTok{1}\NormalTok{]}
\NormalTok{\}}
\NormalTok{padj }\OtherTok{\textless{}{-}} \FunctionTok{p.adjust}\NormalTok{(pvals, }\AttributeTok{method =} \StringTok{"BH"}\NormalTok{)}
\end{Highlighting}
\end{Shaded}

\subsection{4. Check}\label{check}

Volcano plot with the truly DE genes flagged.

\begin{Shaded}
\begin{Highlighting}[]
\FunctionTok{tibble}\NormalTok{(}\AttributeTok{lfc =}\NormalTok{ logfcs, }\AttributeTok{padj =}\NormalTok{ padj, }\AttributeTok{truth =}\NormalTok{ truth) }\SpecialCharTok{|\textgreater{}}
  \FunctionTok{drop\_na}\NormalTok{() }\SpecialCharTok{|\textgreater{}}
  \FunctionTok{ggplot}\NormalTok{(}\FunctionTok{aes}\NormalTok{(lfc, }\SpecialCharTok{{-}}\FunctionTok{log10}\NormalTok{(padj), }\AttributeTok{colour =}\NormalTok{ truth)) }\SpecialCharTok{+}
  \FunctionTok{geom\_point}\NormalTok{(}\AttributeTok{alpha =} \FloatTok{0.6}\NormalTok{) }\SpecialCharTok{+}
  \FunctionTok{labs}\NormalTok{(}\AttributeTok{x =} \StringTok{"log fold change (B vs A)"}\NormalTok{, }\AttributeTok{y =} \StringTok{"{-}log10(padj)"}\NormalTok{)}
\end{Highlighting}
\end{Shaded}

\begin{Shaded}
\begin{Highlighting}[]
\FunctionTok{mean}\NormalTok{(padj[}\DecValTok{11}\SpecialCharTok{:}\NormalTok{n\_gene] }\SpecialCharTok{\textless{}} \FloatTok{0.05}\NormalTok{, }\AttributeTok{na.rm =} \ConstantTok{TRUE}\NormalTok{)}
\end{Highlighting}
\end{Shaded}

\subsection{5. Report}\label{report}

\begin{quote}
On a simulated 500-gene, 10-sample RNA-seq dataset with 10 truly
differentially expressed genes (log fold change ≈ 1.1), per-gene
negative-binomial GLMs followed by Benjamini--Hochberg adjustment
recovered
\texttt{round(mean(padj{[}1:10{]}\ \textless{}\ 0.05,\ na.rm\ =\ TRUE)\ *\ 100,\ 0)}\%
of the true DE genes at FDR \textless{} 0.05 and yielded a
false-discovery rate among null genes of
\texttt{round(mean(padj{[}11:n\_gene{]}\ \textless{}\ 0.05,\ na.rm\ =\ TRUE)\ *\ 100,\ 1)}\%.
\end{quote}

DESeq2 and edgeR would improve power by shrinking gene-wise dispersion
across the experiment; the per-gene fit shown here is a transparent but
low-power version of the same idea.

\subsection{Common pitfalls}\label{common-pitfalls}

\begin{itemize}
\tightlist
\item
  Running a \emph{t}-test on raw counts --- the variance structure is
  wrong.
\item
  Filtering genes after seeing the results (data-driven filters are fine
  only if applied before testing).
\item
  Forgetting to include batch or library-size terms in the design.
\end{itemize}

\subsection{Further reading}\label{further-reading}

\begin{itemize}
\tightlist
\item
  Love MI, Huber W, Anders S (2014), \emph{Moderated estimation of fold
  change and dispersion for RNA-seq data with DESeq2}.
\item
  Law CW et al.~(2014), \emph{voom: precision weights unlock linear
  model analysis tools for RNA-seq read counts}.
\end{itemize}

\subsection{Session info}\label{session-info}

\begin{Shaded}
\begin{Highlighting}[]

\end{Highlighting}
\end{Shaded}

\subsection{See also --- chapter index}\label{see-also-chapter-index}

\begin{itemize}
\tightlist
\item
  \href{../index.Rmd}{Methods in this chapter}
\item
  \href{../../../ch00-find-your-method.Rmd}{Find your method (Chapter
  0)}
\end{itemize}

\end{document}
